% notes.tex - Speaker notes for thesis defense
% Build: xelatex notes.tex
%
% Companion to slides.tex — one page per slide with talking points.

\documentclass[11pt, letterpaper]{article}

\usepackage{fontspec}
\setsansfont{Open Sans}
\newfontfamily\headingfont{Encode Sans}[BoldFont={Encode Sans Bold}]
\renewcommand{\familydefault}{\sfdefault}

\usepackage[margin=0.75in]{geometry}
\usepackage{xcolor}
\usepackage{enumitem}
\usepackage{titlesec}
\usepackage{amssymb}
\usepackage{amsmath}

\definecolor{uwpurple}{RGB}{51, 0, 111}
\definecolor{uwmetallic}{RGB}{145, 123, 76}

% Section = part label
\titleformat{\section}
  {\headingfont\bfseries\Large\color{uwpurple}}
  {}{0em}{}
  [\vspace{-0.3em}{\color{uwmetallic}\rule{3cm}{1.5pt}}]

% Subsection = slide title
\titleformat{\subsection}
  {\headingfont\bfseries\normalsize\color{uwpurple}}
  {Slide \thesubsection:\;}{0em}{}

\setcounter{secnumdepth}{0}
\setcounter{subsection}{0}

% Compact lists
\setlist[itemize]{nosep, left=1em, label=\textcolor{uwpurple}{$\blacktriangleright$}}

\pagestyle{empty}

\begin{document}

{\headingfont\bfseries\LARGE\color{uwpurple} Speaker Notes}\\[2pt]
{\color{uwmetallic}\rule{4cm}{2pt}}\\[6pt]
{\large Cross-Detector Descriptor Fusion --- Frank Sossi, 2026}

\bigskip

% ===== FRAMING =====
\section{Framing}

\stepcounter{subsection}
\subsection{Title Slide}
\begin{itemize}
  \item Welcome everyone. Introduce yourself.
  \item Thank the committee: Prof.\ Olson (chair), Prof.\ Chen, Prof.\ Si.
  \item ``Today I'll present my thesis on cross-detector descriptor fusion for local feature matching.''
\end{itemize}

\stepcounter{subsection}
\subsection{Thesis in One Sentence}
\begin{itemize}
  \item Read the MIOS slowly --- this is the entire thesis in one sentence.
  \item Pause after reading it. This sets up everything that follows.
\end{itemize}

\stepcounter{subsection}
\subsection{Outline}
\begin{itemize}
  \item Four-part structure, about 35--40 minutes then questions.
  \item \textbf{FRAMING} (5 slides):
    \begin{itemize}
      \item The problem --- why local feature matching matters
      \item State of the art --- what existed, what was missing
      \item Key concepts --- detection/description/matching pipeline, mAP
      \item Four research questions with success criteria
      \item Benefits and beneficiaries
    \end{itemize}
  \item \textbf{PROCESS} (9 slides):
    \begin{itemize}
      \item Two evaluation pipelines (full-image vs patch benchmark)
      \item DescriptorWorkbench architecture
      \item Spatial intersection algorithm
      \item Scale characteristics of detectors
      \item Building a color patch benchmark
      \item Same-family fusion fails
      \item Cross-family fusion needs magnitude matching
      \item Pre-fusion L2 normalization solution
      \item Keypoint quality rivals descriptor choice
    \end{itemize}
  \item \textbf{RESULTS} (9 slides):
    \begin{itemize}
      \item Baselines (SIFT 44\%, HardNet 64\%)
      \item Scale control impact (+18\% SIFT, +14\% HardNet)
      \item Validating intersection effect
      \item CNN+CNN fusion to 93.4\%
      \item Discriminator-matcher framework
      \item Patch benchmark fusion results
      \item Answering all four RQs
    \end{itemize}
  \item \textbf{LESSONS LEARNED} (5 slides):
    \begin{itemize}
      \item Technical insights
      \item Engineering and process lessons
      \item Future work
      \item Limitations
      \item Summary of six contributions
    \end{itemize}
\end{itemize}

\stepcounter{subsection}
\subsection{The Problem}
\begin{itemize}
  \item Local features are the backbone of many CV systems.
  \item Detection finds WHERE, description encodes WHAT.
  \item Different families have different strengths --- that's our opportunity.
  \item Key question: can we combine them systematically?
\end{itemize}

\stepcounter{subsection}
\subsection{State of the Art When This Work Began}
\begin{itemize}
  \item LEFT SIDE what existed: 
  \item SIFT is the classic, 128-D gradient histograms from 2004. 
  \item SURF is faster with Haar wavelets, 64-D. 
  \item HardNet and SOSNet are CNN descriptors trained to learn what makes patches similar. 
  \item HoNC and RGBSIFT which capture color infromation that is ignored in the other descriptors.
  \item RIGHT SIDE the gaps: 
  \item There was no work that systematically studied what happens when you fuse descriptors across families. 
  \item There were no good framworks to decide when fusion would help or hurt a system.
  \item While scale effects were noted in the original SIFT paper The role of keypoint scale had not been evaulated
  \item Likewise there had been no work covering detector fusion.
  \item Additionally the original HPatches benchmark only provides grayscale patches, so color descriptors like HoNC and RGBSIfT couldn't be fairly evaluated.
  \item Filling in these gaps define our contribution.
\end{itemize}

\stepcounter{subsection}
\subsection{Key Concepts}
\begin{itemize}
  \item Walk through the pipeline left to right.
  \item \textbf{DETECTION:} find salient locations corners, blobs that are likely to reappear if the camera moves or lighting changes.
  \item \textbf{DESCRIPTION:} at each keypoint, a patch is extracted and encoded as a fixed-length vector. Matching is then done via nearest-neighbor or brute force search in descriptor space.
  \item \textbf{mAP explained:} for each query descriptor, we rank all candidates by distance. Average Precision is the area under the precision-recall curve for that query it rewards ranking the correct match high. Mean AP averages this over all queries across all image pairs. 50\% mAP means on average the correct match is ranked in the top half; 90\% means it's almost always at the top.
\end{itemize}

\stepcounter{subsection}
\subsection{Research Goals \& Measures of Success}
\begin{itemize}
  \item Four research questions, each with a concrete success criterion.
  \item RQ1: detector consensus. RQ2: color fusion. RQ3: compatibility patterns. RQ4: scale impact.
\end{itemize}

\newpage

\stepcounter{subsection}
\subsection{Benefits \& Beneficiaries}
\begin{itemize}
  \item Research community gets: 
  \begin{itemize}
  \item Guidlines for successful fusion.
  \item Evidenc quantifying the effects of keypoint quality.
  \item Verification vs matching framework for classifying descriptors.
  \end{itemize}
  \item Practitioners get: 
  \begin{itemize}
  \item concrete fusion steps and evaultion procedure,
  \item a low cost and easy to implment performance boost from scale filtering, 
  \item An open source tool for evaluating new descriptors.
  \item Applicable to real systems: SLAM, SfM, visual localization all of which use keypoint matching.
  \end{itemize}
\end{itemize}

% ===== PROCESS =====
\newpage
\section{Process}

\stepcounter{subsection}
\subsection{DescriptorWorkbench Architecture}
\begin{itemize}
  \item C++ framework, 10 descriptors, YAML-driven.
  \item SQLite database tracks all 100+ experiments.
  \item Open source on GitHub.
  \item OpenCV 4.13 + LibTorch 2.10
  \item CUDA 13.1 (RTX 4090)
  \item  Google Test for unit tests
  \item Python for analysis
\end{itemize}

\stepcounter{subsection}
\subsection{Two Evaluation Protocols}
\begin{itemize}
  \item \textbf{Full-image pipeline} (Bojanic et al.):
    \begin{itemize}
      \item Image matching (mAP) --- SNN ratio test, 3\,px projection error threshold
      \item Keypoint verification (AUC) --- cross-sequence distractors only
      \item Keypoint retrieval (AP) --- three-tier labeling ($y \in \{-1, 0, +1\}$)
      \item Evaluates detector + descriptor jointly
    \end{itemize}
  \item \textbf{Patch benchmark} (Balntas et al.\ / HPatches):
    \begin{itemize}
      \item Patch matching (mAP) --- rank patches by descriptor distance
      \item Patch verification (AP) --- same-sequence + cross-sequence negatives
      \item Patch retrieval (mAP) --- binary labeling (match vs.\ distractor)
      \item Isolates descriptor quality by holding keypoints constant
    \end{itemize}
  \item Key difference: Bojanic verification uses cross-sequence distractors only; HPatches uses same + cross-sequence negatives.
  \item Both report HP-V (viewpoint) and HP-I (illumination) breakdowns.
\end{itemize}

\stepcounter{subsection}
\subsection{Key Decision: Two Evaluation Pipelines}
\begin{itemize}
  \item This was a \textbf{critical design decision}.
  \item How did we get here?
  \item \textbf{Started with Hpatches} We started by using the hpatches dataset as is but it lacked color information for fusing HonC and RGBSift
  \item \textbf{Locked in Sift Keypoints} We then decided to use locked in keypoints on the original images to get the color information while holding the keypoints constant. When we did this we observed that the performance of non-Sift based descriptors decreased by a considerable margin.
  \item \textbf{Preformance drop Non-sift descriptors} We isolated this effect to differnces in the scale and orientation infromation that each of the non sift detectors had been tuned for. This was especially relavant to the CNN descriptos that were tuned to much higher scales.
  \item \textbf{Implemented intersection} Considering this we implemeted a basic intersection algortithm to find the places that the detectors agreed. The thought behind this was that we could use these keypoints to fairly evaluate each descriptor on the same points in the image but with the scale and orientation information that they were tuned for.
  \item \textbf{Significat perfomance gains} We noticed on these intersection sets the perfomance for all decriptors increased by a significant margin
  \item \textbf{2 pipeines}
  \item Full-image pipeline: Evaluate the use of detector fusion on performance.
  \item Patch pipeline: Re-implement hpatches with color information holds keypoint quality constant isolates descriptor effects.
  \item Without two pipelines, we'd confound our variables.
\end{itemize}

\stepcounter{subsection}
\subsection{Pipline 2 Patch Pipeline: Building a Color Patch Benchmark}
\begin{itemize}
  \item Original HPatches only has grayscale patches.
  \item We \textbf{NEEDED} color for HoNC and RGBSIFT evaluation.
  \item Re-extracted from original images using stored keypoints + homographies.
  \item Validated: SIFT baseline 22.9\% vs 25.47\% original.
  \item Without this, we couldn't answer RQ2 at all.
\end{itemize}

\stepcounter{subsection}
\subsection{Same-Family Fusion Does Not Help}
\begin{itemize}
  \item Patch benchmark isolates descriptor effects.
  \item SIFT alone: 22.9\%, RGBSIFT alone: 24.6\%.
  \item Fusion: 10.4\% --- worse than either alone.
  \item Why? Correlated gradient info. Doubling dimensions adds noise, not signal.
\end{itemize}

\stepcounter{subsection}
\subsection{Cross-Family Fusion Requires Magnitude Matching}
\begin{itemize}
  \item SIFT + HardNet fusion fails without normalization.
  \item Root cause: magnitude mismatch. SIFT values go up to 512, HardNet stays around 0.3.
  \item In L2 distance, SIFT completely dominates --- HardNet contribution is invisible.
\end{itemize}

\stepcounter{subsection}
\subsection{Solution: Pre-Fusion L2 Normalization}
\begin{itemize}
  \item Solution: L2-normalize each component BEFORE fusion. Show the equation.
  \item Without: failed. With: \textbf{46.0\% mAP}.
  \item This led us to make it a configurable option in the framework.
  \item Important engineering lesson: the fix was simple once we found the root cause.
\end{itemize}

\stepcounter{subsection}
\subsection{Spatial Intersection Algorithm}
\begin{itemize}
  \item Walk through the 5 steps. Point to the diagram: blue=SIFT, red=KeyNet, green circles=matches.
  \item Intuition: if two VERY different detectors agree, it's a real feature.
  \item This is a quality filter, not just a count reducer.
  \item \textbf{Algorithm detail --- Mutual Nearest Neighbor (MNN) with spatial tolerance:}
    \begin{itemize}
      \item Build KD-tree spatial indices for both keypoint sets ($T_A$, $T_B$).
      \item For each keypoint $k_a \in K_A$: find nearest neighbor $k_b = \text{NN}(k_a, T_B)$.
      \item Check forward tolerance: $\|k_a - k_b\|_2 \leq \tau$ (Euclidean distance in pixels).
      \item Check mutual agreement: reverse NN of $k_b$ must be $k_a$ (both must be each other's nearest neighbor).
      \item Check uniqueness: each keypoint matched at most once (1-to-1 correspondence).
      \item Output: paired keypoint sets with 1-to-1 correspondence.
    \end{itemize}
  \item \textbf{Only spatial position $(x, y)$ is used for matching} --- scale and orientation are NOT matching criteria. Rationale: SIFT averages 4.45\,px while KeyNet averages 49.83\,px --- requiring scale agreement would eliminate most correspondences.
  \item \textbf{Tolerance $\tau = 3.0$ pixels} (default), based on Mikolajczyk and Schmid's detector evaluation methodology. Strict enough to prevent misalignment, loose enough for typical 1--3\,px detector variance.
  \item \textbf{Tolerance sensitivity:} $\tau=1.0$: 67.41\% mAP (fewest keypoints, highest quality), $\tau=2.0$: 62.10\%, $\tau=5.0$: 57.31\%, $\tau=10.0$: 57.06\% (diminishing returns beyond 5\,px).
  \item \textbf{Example:} SIFT detects corner at $(100.0, 200.0)$, KeyNet detects blob at $(102.1, 201.5)$ --- Euclidean distance 2.65\,px $< \tau=3.0$, so they are paired.
  \item Algorithm is symmetric --- same result regardless of which set is processed first.
\end{itemize}

\stepcounter{subsection}
\subsection{Full Image Pipeline: Scale Characteristics of Detectors}
\begin{itemize}
  \item SIFT: tiny keypoints, average \textbf{4.45 pixels}.
  \item KeyNet: much larger, average \textbf{49.83 pixels}.
  \item Key insight: bigger keypoint = bigger patch = more information.
  \item 4px samples 16$\times$16, 10px samples 40$\times$40 --- huge difference.
\end{itemize}

\stepcounter{subsection}
\subsection{Keypoint Quality Matters More Than Descriptor Choice}
\begin{itemize}
  \item Descriptor research focuses on encoding algorithms. But keypoint selection matters just as much.
  \item As shown in the table we get comparable gains in mAP by using intersections sets as the gain achieved by using the more complicated CNN descriptors.
  \item These are comparable gains.
\end{itemize}

% ===== RESULTS =====
\newpage
\section{Results}

\stepcounter{subsection}
\subsection{Baseline Performance}
\begin{itemize}
  \item Starting point for all comparisons.
  \item SIFT $\sim$44\%, RootSIFT $\sim$47\%, HardNet/SOSNet $\sim$64\%.
  \item CNN descriptors $\sim$20\% better than traditional.
  \item SIFT prefers viewpoint, CNN prefers illumination.
  \item All improvements in next slides measured from these baselines.
\end{itemize}

\stepcounter{subsection}
\subsection{Result: Scale Control Impact}
\begin{itemize}
  \item Point to the figure.
  \item SIFT: 44.5\% to \textbf{62.8\%} = +18.3\% absolute.
  \item HardNet: 64.5\% to \textbf{78.1\%} = +13.6\% absolute.
  \item Just filtering to top 25\% by scale!
  \item Quality over quantity --- fewer but better keypoints.
  \item Answers RQ4: scale has large, consistent impact.
\end{itemize}

\stepcounter{subsection}
\subsection{Result: Detector Intersection Progression}
\begin{itemize}
  \item Three stages shown in the figure.
  \item Full set: 64.5\% $\to$ Scale-controlled: 78.1\% $\to$ Intersection: \textbf{82.1\%}.
  \item Each stage adds cumulative improvement.
  \item Intersection adds +4\% BEYOND scale control.
  \item Answers RQ1: yes, detector consensus is a quality signal.
\end{itemize}

\stepcounter{subsection}
\subsection{Validating the Intersection Mechanism (1/2)}
\begin{itemize}
  \item The +18\% gain required verification. We ruled out four alternatives.
  \item NOT better repeatability: Intersection 28.0\% vs Scale 29.4\%.
  \item NOT more distinctive descriptors: NN ratios identical ($\sim$0.44).
  \item NOT higher response values: both sets avg $\sim$0.035.
  \item NOT spatial crowding: quadrant distributions identical.
\end{itemize}

\stepcounter{subsection}
\subsection{Validating the Intersection Mechanism (2/2)}
\begin{itemize}
  \item Confirmed: detector consensus removes ``confusing'' keypoints at repetitive textures.
  \item Higher precision at threshold 0.8: 71.2\% vs 66.2\%.
  \item Fewer false positives: 0.40 vs 0.51 FP per TP (22\% reduction).
  \item Key insight: if ALL descriptors benefit equally, the locations are better, not the descriptors.
\end{itemize}

\stepcounter{subsection}
\subsection{Result: CNN + CNN Fusion}
\begin{itemize}
  \item On intersection keypoints.
  \item HardNet alone: 82.1\%, SOSNet alone: 82.0\%.
  \item Concatenation: \textbf{93.4\%} --- that's +11.3\%!
  \item Averaging: 92.3\% --- good but concatenation is better.
  \item They learn complementary representations despite similar training.
\end{itemize}

\stepcounter{subsection}
\subsection{The Discriminator--Matcher Framework}
\begin{itemize}
  \item V/M ratio = verification / matching performance.
  \item HoNC: \textbf{3.84$\times$} --- great at rejecting false matches (discriminator).
  \item HardNet: \textbf{1.84$\times$} --- trained for correspondence (matcher).
  \item This framework PREDICTS fusion outcomes.
  \item Pairing discriminator + matcher should be best.
\end{itemize}

\stepcounter{subsection}
\subsection{Result: Patch Benchmark Fusion}
\begin{itemize}
  \item Point to the figure.
  \item Best: HoNC + SOSNet = \textbf{50.6\%} (discriminator + matcher).
  \item CNN + CNN = 49.9\% (complementary learned features).
  \item SIFT + HoNC = 15.5\% --- WORSE than either alone (two weak matchers).
  \item Answers RQ2 and RQ3: color adds value with strong matcher, complementarity is key.
\end{itemize}

\stepcounter{subsection}
\subsection{Answering the Research Questions}
\begin{itemize}
  \item Go through each RQ with its answer.
  \item RQ1: YES --- intersection +18\%.
  \item RQ2: YES --- HoNC+SOSNet beats SOSNet alone.
  \item RQ3: complementarity determines success, magnitude matching required.
  \item RQ4: scale is dominant --- comparable to switching descriptor families.
  \item All four questions answered with quantitative evidence.
\end{itemize}

% ===== LESSONS LEARNED =====
\section{Lessons Learned}

\stepcounter{subsection}
\subsection{What I Learned: Technical Insights}
\begin{itemize}
  \item Three key lessons.
  \item 1.\ Keypoint quality is \textbf{underappreciated} in the literature.
  \item 2.\ Investigating failures led to the magnitude matching discovery.
  \item 3.\ Two pipelines prevented false conclusions about fusion.
\end{itemize}

\stepcounter{subsection}
\subsection{What I Learned: Engineering \& Process}
\begin{itemize}
  \item 1.\ Experiment infrastructure (SQLite + YAML) was worth the investment.
  \item 2.\ Reproducibility requires tooling --- every result from one YAML file.
  \item 3.\ The V/M ratio framework emerged from data visualization, not theory.
\end{itemize}

\stepcounter{subsection}
\subsection{Future Work}
\begin{itemize}
  \item Short-term: learned weights, tolerance sensitivity, more descriptors.
  \item Long-term: end-to-end learning, other datasets, real-time deployment.
  \item Plenty of directions to extend this work.
\end{itemize}

\stepcounter{subsection}
\subsection{Limitations}
\begin{itemize}
  \item Single dataset (HPatches) --- may not generalize.
  \item Concatenation doubles dimensionality (but filtering reduces total cost 28$\times$).
  \item Fixed equal weighting --- learned weights could help.
\end{itemize}

\stepcounter{subsection}
\subsection{Summary of Contributions}
\begin{itemize}
  \item Read through the six contributions.
  \item Emphasize the bottom-line message: \textbf{keypoint selection matters as much as descriptor choice}.
  \item This is the take-home message.
\end{itemize}

\stepcounter{subsection}
\subsection{Thank You / Questions}
\begin{itemize}
  \item Thank the committee again. Thank the audience.
  \item Open for questions.
  \item GitHub link available for anyone interested.
\end{itemize}

\end{document}
